\chapter{Conclusion}
\label{chap:conclusion}

% =====================================================================
%  RAPPEL INSTRUCTION : un bilan FONCTIONNEL (ce qui était demandé a-t-il
%  été fait ? les tests ont-ils été réalisés ?) ET une perspective sur ce
%  que l'alternance t'a apporté (technique, managérial, personnel).
% =====================================================================

\section{Bilan fonctionnel}
Au moment de la rédaction de ce mémoire, l'alternance n'est pas achevée :
les missions menées chez Beemo Technologie sont encore en cours. Le bilan
qui suit photographie donc l'état des travaux à la date de la soutenance.

Côté Free Mobile, les cinq missions ont été menées à leur terme et livrées.
L'outil de duplication Grafana (chapitre~\ref{chap:t1}) est exploité pour
synchroniser les instances de supervision. Le framework de test
(chapitre~\ref{chap:t3}) a remplacé la suite Lua sur l'ensemble de son
périmètre et constitue désormais le filet de non-régression du \emph{Valo
Engine} ; l'extension VS~Code (chapitre~\ref{chap:t4}) l'accompagne et en a
ouvert l'usage à toute l'équipe. Le portage vers les forks legacy
(chapitre~\ref{chap:t5}) a couvert les six variantes visées — France B2C et
B2B, Caraïbes, Italie B2C et B2B, proxy italien — et a été l'occasion d'une
transmission de connaissances en binôme avant mon départ. Enfin, l'étude sur
le coût de compilation (chapitre~\ref{chap:t2}) a rendu le \emph{build}
soutenable sur les machines de développement, tout en documentant la dette de
fond que seuls les modules C++20 résorberont à terme.

Côté Beemo, la rétro-ingénierie du DSDK (chapitre~\ref{chap:t6}) a atteint
son objectif principal — un \emph{build} local, reproductible et versionné —
mais reste ouverte sur la validation du package affecté par l'un des
correctifs. La fonctionnalité d'immuabilité (chapitre~\ref{chap:t7}) est en
cours de développement : ses briques centrales — automate à états, service de
chiffrement, couche d'adaptation au firmware — sont en place, le code neuf
étant écrit de manière testable et isolée, à rebours des pratiques du
firmware qui l'accueille.

\section{Retour sur le fil rouge : la dette technique}
Relues à travers la grille du chapitre~\ref{chap:veille}, ces sept missions
dessinent une conclusion d'ensemble : la dette technique ne se rembourse pas
par un grand soir, mais par une accumulation d'actions ciblées, chacune
actionnant un levier précis — outiller, tester, fiabiliser la chaîne de
production, isoler le neuf de l'ancien. Dans les deux entreprises, la
stratégie gagnante a été \emph{incrémentale} : préserver la capacité à livrer
tout en réduisant, intervention après intervention, le risque et le coût des
suivantes.

J'en retiens aussi deux enseignements moins attendus. Le premier est que la
forme de dette la plus dangereuse n'est pas dans le code : c'est la dette de
connaissance, celle qui part avec les personnes — le framework de test après
le départ de David, la procédure de compilation de l'agent VSX que je suis
venu reprendre. Sa parade est la capitalisation sous forme exécutable :
tests, patchs versionnés, outillage. Le second est la dimension humaine du
remboursement : un travail de fond sur du legacy est ingrat, et la motivation
des équipes en est un facteur limitant au même titre que la technique. Les
outils d'\gls{ia-agentique}, expérimentés en fin de parcours
(chapitre~\ref{chap:t5}), ne changent pas cette équation : ils récompensent
la maîtrise de l'architecture et du métier, mais ne s'y substituent pas.

\section{Apports de l'alternance}
Sur le plan technique, ces deux années m'ont fait pratiquer un spectre
inhabituellement large pour un alternant : le C++ d'un moteur critique et
d'un firmware embarqué, le Python d'un framework de test de 30~000 lignes
conçu et architecturé de bout en bout, le TypeScript d'une extension
VS~Code, les protocoles \gls{lsp} et \gls{dap}, la rétro-ingénierie d'un SDK
propriétaire, Docker, et les arcanes d'une chaîne de compilation C++. Plus
que les technologies elles-mêmes, j'en retire une compétence transversale :
savoir aborder un système existant, opaque et non documenté, et en reprendre
le contrôle méthodiquement.

Sur le plan organisationnel et humain, travailler dans deux entreprises —
un grand opérateur et une PME d'une trentaine de personnes — m'a donné deux
points de comparaison sur la méthode agile, la gestion d'une équipe face au
legacy, et la place qu'un développeur peut y occuper : transverse et hors
sprint chez Free Mobile, pleinement intégré aux livraisons chez Beemo. Le
portage en binôme m'a par ailleurs placé en situation de \emph{transmettre} —
exercice qui force à expliciter ce que l'on croyait évident. Enfin, sur le
plan personnel, la décision de quitter Free Mobile pour suivre une nouvelle
opportunité m'a appris à arbitrer un choix de carrière, et la confiance
accordée sur des missions structurantes a confirmé mon goût pour
l'ingénierie au service de la qualité logicielle.

\section{Perspectives}
Plusieurs chantiers prolongent naturellement ces travaux. Chez Free Mobile,
l'architecture \emph{core}/\emph{spec} du framework prépare la modernisation
des forks legacy vers un socle commun, le filet de non-régression rendant
désormais cette migration envisageable ; la résorption de la dette de
compilation passera, elle, par l'adoption des modules C++20 lorsque les
chaînes de compilation l'autoriseront. Chez Beemo, la suite est double :
fiabiliser la validation des packages auprès de Broadcom — puis l'automatiser,
pour achever l'harmonisation de la chaîne de production de l'agent VSX — et
mener à son terme la fonctionnalité d'immuabilité. À plus long terme, la
couche d'adaptation qui l'héberge ouvre une voie de modernisation progressive
du firmware : écrire chaque évolution en C++ moderne et testable, et « gommer »
le legacy fonctionnalité par fonctionnalité.
